\ifdefined\gkver\else\def\gkver{student}\fi
\documentclass[\gkver]{gaokaozhenti}
\begin{document}

\chapter{2005年上海}

\begin{enumerate}
\item
%% number: 1
%% paperName: 2005年上海高考第 1 题 4 分
%% typeId: 3
%% type: 填空题
%% chapter: 磁场
%% point: 左手定则
%% method: 无
%% score: 4
%% degree: 850
%% duplicateId: 0
%% body:
\textbf{A．}通电直导线 A 与圆形通电导线环 B 固定放置在同一水平面上，通有如图所示的电流时，通电直导线 A 受到水平向\tkanswer[1.2cm]{右}的安培力作用。当 A、B 中电流大小保持不变，但同时改变方向时，通电直导线 A 所受到的安培力方向水平向\tkanswer[1.2cm]{右}。

\onepicture[w=3.2cm, align=right]{figs/01a.png}

\textbf{B．}右面是逻辑电路图及其真值表，此逻辑电路为\tkanswer[1.2cm]{“或”}门电路，在真值表中 X 处的逻辑值为\tkanswer[1.2cm]{1}。

\begin{center}
\begin{tabular}{|c|c|c|}
\hline
\multicolumn{2}{|c|}{输入} & 输出 \\ \hline
A & B & Z \\ \hline
0 & 0 & 0 \\ \hline
0 & 1 & 1 \\ \hline
1 & 0 & X \\ \hline
1 & 1 & 1 \\ \hline
\end{tabular}
\end{center}

\onepicture[w=3.4cm, align=right]{\includesvg{figs/01b}}

%% answer:
\jdanswer{
\textbf{A．}右，右。

\textbf{B．}“或”，1。
}
%% memo:
\memoanswer{
\textbf{A．}由安培定则可得，A 所在处的磁场垂直纸面向里，则由左手定则可得，A 受到的安培力向右；当 A、B 中电流同时反向后，由安培定则可得，A 所在处的磁场垂直纸面向外，则由左手定则可得，A 受到的安培力仍向右。

\textbf{B．}逻辑电路分为“或”门电路、“与”门电路和“非”门电路。由图表中输入和输出的关系，可知本题应为“或”门电路，得出 X 值应为 1。
}

\item
%% number: 2
%% paperName: 2005年上海高考第 2 题 4 分
%% typeId: 3
%% type: 填空题
%% chapter: 交流电
%% point: 交流电
%% method: 无
%% score: 4
%% degree: 850
%% duplicateId: 0
%% body:
\textbf{A．}如图所示，实线表示两个相干波源 $S_{1}$、$S_{2}$ 发出的波的波峰位置，则图中的\tkanswer[1.0cm]{b}点为振动加强的位置，图中的\tkanswer[1.0cm]{a}点为振动减弱的位置。

\onepicture[w=4.6cm, align=right]{\includesvg{figs/02a}}

\textbf{B．}正弦交流电是由闭合线圈在匀强磁场中匀速转动产生的。线圈中感应电动势随时间变化的规律如图所示，则此感应电动势的有效值为\tkanswer[2.0cm]{$220$（或 $\frac{311\sqrt{2}}{2}$）}$\UV$，频率为\tkanswer[1.2cm]{50}$\UHz$。

\onepicture[w=5.0cm, align=right]{\includesvg{figs/02b}}

%% answer:
\jdanswer{
\textbf{A．}b，a。

\textbf{B．}220（或 $\frac{311\sqrt{2}}{2}$），50。
}
%% memo:
\memoanswer{
\textbf{A．}在波的传播过程中波峰和波谷交替出现，故在两个相邻波峰之间一定存在一个波谷，故 b 点应该既是 $S_{1}$ 的波谷位置，又是 $S_{2}$ 的波谷位置，即为两个波谷的重叠区域，故 b 点是振动加强位置；而 a 点是 $S_{2}$ 的波峰位置，也是 $S_{1}$ 的波谷位置，故是振动减弱的位置。

\textbf{B．}由 $e-t$ 图象可知，感应电动势的最大值为 $311\UV$，故有效值为 $E=\frac{E_{m}}{\sqrt{2}}=\frac{311}{\sqrt{2}}\UV=220\UV$；由 $e-t$ 图象可知，交变电流的周期 $T=0.02\Us$，所以频率为 $f=\frac{1}{T}=50\UHz$。
}

\item
%% number: 3
%% paperName: 2005年上海高考第 3 题 4 分
%% typeId: 3
%% type: 填空题
%% chapter: 直线运动
%% point: 自由落体运动
%% method: 无
%% score: 4
%% degree: 850
%% duplicateId: 0
%% body:
\textbf{A．}对“落体运动快慢”“力与物体运动关系”等问题，亚里士多德和伽利略存在着不同的观点。请完成下表：

\begin{center}
\begin{tabular}{|c|c|c|}
\hline
 & 亚里士多德的观点 & 伽利略的观点 \\ \hline
落体运动快慢 & 重的物体下落快，轻的物体下落慢 & \tkanswer[3.2cm]{物体下落快慢与物体轻重无关} \\ \hline
力与物体运动关系 & \tkanswer[2.6cm]{维持物体运动需要力} & 维持物体运动不需要力 \\ \hline
\end{tabular}
\end{center}

\textbf{B．}阴极射线是从阴极射线管的阴极发出的高速运动的粒子流，这些微观粒子是\tkanswer[1.2cm]{电子}。若在如图所示的阴极射线管中部加上垂直于纸面向里的磁场，阴极射线将\tkanswer[1.2cm]{向下}（填“向上”“向下”“向里”“向外”）偏转。

\onepicture[w=5.6cm, align=right]{\includesvg{figs/03b}}

%% answer:
\jdanswer{
\textbf{A．}物体下落快慢与物体轻重无关，维持物体运动需要力。

\textbf{B．}电子，向下。
}
%% memo:
\memoanswer{
\textbf{A．}亚里士多德根据日常生活经验，不通过实验验证认为重的物体下落快、轻的物体下落慢，维持物体运动需要力；伽利略通过对自由落体运动的研究，推翻了他的观点，伽利略认为物体下落快慢与物体轻重无关，维持物体运动不需要力。

\textbf{B．}阴极射线是金属丝受热后产生的高速电子流；电子定向移动的过程中形成了电流，电子从阴极向阳极移动，可等效为在阴阳极之间产生了从右至左的电流，由左手定则知，电流受到向下的安培力，因此射线向下偏转。
}

\item
%% number: 4
%% paperName: 2005年上海高考第 4 题 4 分
%% typeId: 3
%% type: 填空题
%% chapter: 电场
%% point: 电场叠加
%% method: 对称法
%% score: 4
%% degree: 650
%% duplicateId: 0
%% body:
如图，带电量为 $+q$ 的点电荷与均匀带电薄板相距为 $2d$，点电荷到带电薄板的垂线通过板的几何中心。若图中 a 点处的电场强度为零，根据对称性，带电薄板在图中 b 点处产生的电场强度大小为\tkanswer[2.0cm]{$\frac{kq}{d^{2}}$}，方向\tkanswer[3.4cm]{水平向左（或垂直薄板向左）}。（静电力恒量为 $k$）

\onepicture[w=4.6cm, align=right]{\includesvg{figs/04}}

%% answer:
\jdanswer{$\frac{kq}{d^{2}}$，水平向左（或垂直薄板向左）}
%% memo:
\memoanswer{
由于 a 点的合电场强度为零，则点电荷 $+q$ 在 a 点产生的场强 $E_{1}$ 和带电薄板在 a 点产生的场强 $E_{a}$ 大小相等，即 $E_{1}=E_{a}$，方向相反；由于 a、b 两点相对薄板对称，所以带电薄板在 b 点产生的场强 $E_{b}$ 和在 a 点产生的场强 $E_{a}$ 等大反向，即 $E_{b}=E_{1}=\frac{kq}{d^{2}}$，方向和点电荷 $q$ 产生的场强 $E_{1}$ 方向相同，故正确答案为 $\frac{kq}{d^{2}}$，方向水平向左（或垂直薄板向左）。
}

\item
%% number: 5
%% paperName: 2005年上海高考第 5 题 4 分
%% typeId: 3
%% type: 填空题
%% chapter: 电路
%% point: 电路中的图象
%% method: 无
%% score: 4
%% degree: 650
%% duplicateId: 0
%% body:
右图中图线①表示某电池组的输出电压—电流关系，图线②表示其输出功率—电流关系。该电池组的内阻为\tkanswer[1.2cm]{5}$\UO$。当电池组的输出功率为 $120\UW$ 时，电池组的输出电压是\tkanswer[1.2cm]{30}$\UV$。

\onepicture[w=7.0cm, align=right]{\includesvg{figs/05}}

%% answer:
\jdanswer{5，30}
%% memo:
\memoanswer{
图线①是输出电压和电流的关系图象，当输出电流为 0 时，输出电压即为电源电动势，可得电源的电动势为 $50\UV$；当输出电压为 $40\UV$ 时，输出电流为 $2\UA$，所以电源内阻两端电压为 $10\UV$，由欧姆定律可得电源的内阻为 $5\UO$。由图线②可知电源输出功率为 $120\UW$ 时，输出电流为 $4\UA$，由 $P=UI$ 可得电池的输出电压为 $30\UV$。
}

\item
%% number: 6
%% paperName: 2005年上海高考第 6 题 5 分
%% typeId: 2
%% type: 多选题
%% chapter: 光学
%% point: 光电效应
%% method: 无
%% score: 5
%% degree: 650
%% duplicateId: 0
%% body:
2005 年被联合国定为“世界物理年”，以表彰爱因斯坦对科学的贡献。爱因斯坦对物理学的贡献有\xzanswer{AC}
\fourchoices[answer=AC]
{创立“相对论”}
{发现“X 射线”}
{提出“光子说”}
{建立“原子核式模型”}
%% answer: AC
%% memo:
\memoanswer{
爱因斯坦创立了“相对论”和提出了“光子说”；伦琴发现了“X 射线”，卢瑟福建立了“原子核式模型”，故 AC 正确。
}

\item
%% number: 7
%% paperName: 2005年上海高考第 7 题 5 分
%% typeId: 2
%% type: 多选题
%% chapter: 原子和原子核
%% point: 人工核反应
%% method: 无
%% score: 5
%% degree: 850
%% duplicateId: 0
%% body:
卢瑟福通过实验首次实现了原子核的人工转变，核反应方程为 $\ce{^{4}_{2}He + ^{14}_{7}N -> ^{17}_{8}O + ^{1}_{1}H}$，下列说法中正确的是\xzanswer{AC}
\fourchoices[answer=AC]
{通过此实验发现了质子}
{实验中利用了放射源放出的 $\gamma$ 射线}
{实验中利用了放射源放出的 $\alpha$ 射线}
{原子核在人工转变过程中，电荷数可能不守恒}
%% answer: AC
%% memo:
\memoanswer{
卢瑟福利用放射源放出的 $\alpha$ 射线通过该实验发现了质子，故 AC 正确、B 错误；核反应满足质量数和电荷数守恒，故 D 错误。
}

\item
%% number: 8
%% paperName: 2005年上海高考第 8 题 5 分
%% typeId: 2
%% type: 多选题
%% chapter: 力物体的平衡
%% point: 摩擦力
%% method: 无
%% score: 5
%% degree: 650
%% duplicateId: 0
%% body:
对如图所示的皮带传动装置，下列说法中正确的是\xzanswer{BD}
\onepicture[w=4.2cm]{\includesvg{figs/08}}
\fourchoices[answer=BD]
{A 轮带动 B 轮沿逆时针方向旋转}
{B 轮带动 A 轮沿逆时针方向旋转}
{C 轮带动 D 轮沿顺时针方向旋转}
{D 轮带动 C 轮沿顺时针方向旋转}
%% answer: BD
%% memo:
\memoanswer{
对于 AB，由图可以看出下方的皮带是被拉紧的，上方的皮带没有承受拉力，所以如果 A 是主动轮它将顺时针旋转，如果 B 是主动轮则它要逆时针旋转；CD 中上面皮带被拉紧，分析方法类似，故 BD 正确。
}

\item
%% number: 9
%% paperName: 2005年上海高考第 9 题 5 分
%% typeId: 2
%% type: 多选题
%% chapter: 功和能
%% point: 机械能守恒定律
%% method: 无
%% score: 5
%% degree: 650
%% duplicateId: 0
%% body:
如图所示，A、B 分别为单摆做简谐振动时摆球的不同位置。其中，位置 A 为摆球摆动的最高位置，虚线为过悬点的竖直线。以摆球最低位置为重力势能零点，则摆球在摆动过程中\xzanswer{BC}
\onepicture[w=4.0cm]{\includesvg{figs/09}}
\fourchoices[answer=BC]
{位于 B 处时动能最大}
{位于 A 处时势能最大}
{在位置 A 的势能大于在位置 B 的动能}
{在位置 B 的机械能大于在位置 A 的机械能}
%% answer: BC
%% memo:
\memoanswer{
在单摆摆动的过程中，小球的机械能守恒，小球经过最低点时重力势能最小，动能最大，小球经过最高点（如 A 点）时动能最小，势能最大。由于 B 不是最高点，所以 B 的速度不是 0，动能不是 0，由于机械能守恒，其动能小于 A 点的重力势能，故 BC 正确。
}

\item
%% number: 10
%% paperName: 2005年上海高考第 10 题 5 分
%% typeId: 2
%% type: 多选题
%% chapter: 曲线运动
%% point: 类平抛运动
%% method: 无
%% score: 5
%% degree: 650
%% duplicateId: 0
%% body:
如图所示的塔吊臂上有一可以沿水平方向运动的小车 A，小车下装有吊着物体 B 的吊钩。在小车 A 与物体 B 以相同的水平速度沿吊臂方向匀速运动的同时，吊钩将物体 B 向上吊起，A、B 之间的距离以 $d=H-2t^{2}$（SI）（SI 表示国际单位制，式中 $H$ 为吊臂离地面的高度）规律变化，则物体做\xzanswer{BC}
\onepicture[w=4.0cm]{\includesvg{figs/10}}
\fourchoices[answer=BC]
{速度大小不变的曲线运动}
{速度大小增加的曲线运动}
{加速度大小方向均不变的曲线运动}
{加速度大小方向均变化的曲线运动}
%% answer: BC
%% memo:
\memoanswer{
由于 B 在水平方向上做匀速直线运动，在竖直方向上由公式 $d=H-2t^{2}$ 可知其做初速度为 0 的匀加速直线运动，这两个方向的运动合成后便是匀变速曲线运动，故 BC 正确。
}

\item
%% number: 11
%% paperName: 2005年上海高考第 11 题 5 分
%% typeId: 1
%% type: 单选题
%% chapter: 电磁感应
%% point: 电磁感应中图象
%% method: 无
%% score: 5
%% degree: 650
%% duplicateId: 0
%% body:
如图所示，A 是长直密绕通电螺线管。小线圈 B 与电流表连接，并沿 A 的轴线 $Ox$ 从 $O$ 点自左向右匀速穿过螺线管 A。能正确反映通过电流表中电流 $I$ 随 $x$ 变化规律的是\xzanswer{C}
\onepicture[w=5.6cm]{\includesvg{figs/11}}
\fourchoices[answer=C, ispicture=true, h=2.0cm]
{figs/11a.png}
{figs/11b.png}
{figs/11c.png}
{figs/11d.png}
%% answer: C
%% memo:
\memoanswer{
通电螺线管内部产生的是匀强磁场，外部的磁场和条形磁体相似，故 B 在从 O 点进入螺线管时通过 B 的磁通量是增加的，到中间后由于进入了匀强磁场，通过 B 的磁通量不再变化，故没有了感应电流；当 B 从内部出来时通过 B 的磁通量是在减小的，所以在 B 中会产生一个和进入时方向相反的感应电流，故 C 正确。
}

\item
%% number: 12
%% paperName: 2005年上海高考第 12 题 5 分
%% typeId: 2
%% type: 多选题
%% chapter: 电场
%% point: 带电粒子的直线运动
%% method: 无
%% score: 5
%% degree: 650
%% duplicateId: 0
%% body:
在场强大小为 $E$ 的匀强电场中，一质量为 $m$、带电量为 $q$ 的物体以某一初速沿电场反方向做匀减速直线运动，其加速度大小为 $0.8qE/m$，物体运动 $s$ 距离时速度变为零。则\xzanswer{ACD}
\fourchoices[answer=ACD]
{物体克服电场力做功 $qEs$}
{物体的电势能减少了 $0.8qEs$}
{物体的电势能增加了 $qEs$}
{物体的动能减少了 $0.8qEs$}
%% answer: ACD
%% memo:
\memoanswer{
带电量为 $+q$ 的物体在场强大小为 $E$ 的匀强电场中所受的电场力大小为 $qE$，方向沿电场线方向。当物体逆电场线方向的位移为 $s$ 时电场力做负功 $qEs$，电势能增加 $qEs$；对于物体动能的变化可以由动能定理来判断，由物体的加速度可得物体所受的合外力为 $0.8qE$，故发生位移 $s$ 时其动能减小 $0.8qEs$，故 ACD 正确。
}

\item
%% number: 13
%% paperName: 2005年上海高考第 13 题 5 分
%% typeId: 2
%% type: 多选题
%% chapter: 机械波
%% point: 波的多解性
%% method: 无
%% score: 5
%% degree: 450
%% duplicateId: 0
%% body:
A、B 两列波在某时刻的波形如图所示，经过 $t=T_{\mathrm{A}}$ 时间（$T_{\mathrm{A}}$ 为波 A 的周期），两波再次出现如图波形，则两波的波速之比 $v_{\mathrm{A}}:v_{\mathrm{B}}$ 可能是\xzanswer{ABC}
\twopicture[wA=6.0cm, wB=6.0cm]{figs/13a.png}{figs/13b.png}
\fourchoices[answer=ABC]
{$1:3$}
{$1:2$}
{$2:1$}
{$3:1$}
%% answer: ABC
%% memo:
\memoanswer{
已知波速公式为 $v=\frac{\lambda}{T}$，当经过时间 $t$ 后 B 再次出现如图所示的波形，则 $t=nT_{B}$（$n$ 为正整数），即 $T_{A}=nT_{B}$；由图可知 $\lambda_{A}=2\lambda_{B}$，则 $v_{B}=\frac{\lambda_{B}}{T_{B}}=\frac{n\lambda_{A}}{2T_{A}}=\frac{n}{2}v_{A}$，即 $\frac{v_{A}}{v_{B}}=\frac{2}{n}$，当 $n$ 取不同值时，可知 ABC 正确。
}

\item
%% number: 14
%% paperName: 2005年上海高考第 14 题 6 分
%% typeId: 3
%% type: 填空题
%% chapter: 光学
%% point: 光的衍射
%% method: 无
%% score: 6
%% degree: 850
%% duplicateId: 0
%% body:
部分电磁波的大致波长范围如图所示。若要利用缝宽与手指宽度相当的缝获得明显的衍射现象，可选用\tkanswer[1.2cm]{微波}波段的电磁波，其原因是\tkanswer[4.2cm]{要产生明显的衍射现象，波长应与缝的尺寸相近}。

\onepicture[w=7.0cm, align=right]{\includesvg{figs/14}}

%% answer:
\jdanswer{微波；要产生明显的衍射现象，波长应与缝的尺寸相近}
%% memo:
\memoanswer{
微波；要产生明显的衍射现象，波长应与缝的尺寸相近。
}

\item
%% number: 15
%% paperName: 2005年上海高考第 15 题 6 分
%% typeId: 4
%% type: 实验题
%% chapter: 气体性质
%% point: 玻意耳定律
%% method: 无
%% score: 6
%% degree: 450
%% duplicateId: 0
%% body:
一同学用下图装置研究一定质量气体的压强与体积的关系，实验过程中温度保持不变。最初，U 形管两臂中的水银面齐平，烧瓶中无水。当用注射器往烧瓶中注入水时，U 形管两臂中的水银面出现高度差。实验的部分数据记录在右表。

\begin{center}
\begin{tabular}{|c|c|c|c|c|c|}
\hline
气体体积 $V$（$\UmL$） & 800 & 674 & 600 & 531 & 500 \\ \hline
水银面高度差 $h$（$\Ucm$） & 0 & 14.0 & 25.0 & 38.0 & 45.0 \\ \hline
\end{tabular}
\end{center}

\begin{enumerate}
	\item
	根据表中数据，在右图中画出该实验的 $h-1/V$ 关系图线。
	\item
	实验时，大气压强 $p_{0}=$\tkanswer[2.0cm]{75.0}$\UcmHg$。
\end{enumerate}

\onepicture[h=4.6cm, align=right]{\includesvg{figs/15a}}
\onepicture[h=6.0cm, align=right]{\includesvg{figs/15b}}

%% answer:
\jdanswer{
\begin{enumerate}
	\item
	如图所示。
	\item
	$75.0\UcmHg$（$74.5\UcmHg\sim75.5\UcmHg$）。
\end{enumerate}
}
%% memo:
\memoanswer{
（1）气体压强与体积的关系图像如图所示。

\onepicture[h=6.0cm]{\includesvg{figs/15_an}}

（2）设大气压强为 $p_{0}$，由压强平衡得 $p=p_{0}+h$，由气体实验定律得 $\frac{pV}{T}=C$（常数），联立化简得 $h=\frac{CT}{V}-p_{0}$，故图像纵截距的绝对值为大气压强，由图像知 $p_{0}=75\UcmHg$。
}

\item
%% number: 16
%% paperName: 2005年上海高考第 16 题 6 分
%% typeId: 4
%% type: 实验题
%% chapter: 电路
%% point: 分压与分流
%% method: 无
%% score: 6
%% degree: 650
%% duplicateId: 0
%% body:
一根长为 $1\Um$ 的均匀电阻丝需与一“$10\UV$，$5\UW$”的灯同时工作，电源电压恒为 $100\UV$，电阻丝阻值 $R=100\UO$（其阻值不随温度变化）。现利用分压电路从电阻丝上获取电能，使灯正常工作。

\begin{enumerate}
	\item
	在右面方框中完成所需电路；
	\item
	电路中电流表的量程应选择\tkanswer[2.4cm]{$0\sim 3\UA$}（选填：“$0\sim 0.6\UA$”或“$0\sim 3\UA$”）；
	\item
	灯正常工作时，与其并联的电阻丝长度为\tkanswer[1.6cm]{0.17}$\Um$（计算时保留小数点后二位）。
\end{enumerate}

\onepicture[w=5.2cm, align=right]{\includesvg{figs/16}}

%% answer:
\jdanswer{
\begin{enumerate}
	\item
	如图所示。
	\item
	$0\sim 3\UA$。
	\item
	$0.17$（或 $\frac{3\sqrt{5}-5}{10}$）。
\end{enumerate}
}
%% memo:
\memoanswer{
（1）利用分压接法应先将电阻丝接入电路，再将灯泡与电阻丝的一部分并联，电流表接在干路中，如图所示。

\onepicture[w=5.2cm]{\includesvg{figs/16_an}}

（2）因电路中电流最大为 $\frac{100\UV}{100\UO}=1\UA$，为了保证安全应选 $0\sim 3\UA$ 的量程。

（3）灯泡正常工作，则灯泡两端的电压应为 $10\UV$，灯泡的电阻 $R=\frac{U^{2}}{P}=20\UO$。设并联的电阻丝长度为 $x$，则并联的电阻应为 $100x$，由串联电路的电压分配的特点知 $\frac{\frac{100x\times20}{100x+20}}{100(1-x)}=\frac{1}{9}$，解得 $x=\frac{3\sqrt{5}-5}{10}\Um\approx0.17\Um$。
}

\item
%% number: 17
%% paperName: 2005年上海高考第 17 题 7 分
%% typeId: 4
%% type: 实验题
%% chapter: 电路
%% point: 故障分析
%% method: 无
%% score: 7
%% degree: 450
%% duplicateId: 0
%% body:
两实验小组使用相同规格的元件，按右图电路进行测量。他们将滑动变阻器的滑片 P 分别置于 a、b、c、d、e 五个间距相同的位置（a、e 为滑动变阻器的两个端点），把相应的电流表示数记录在表一、表二中。对比两组数据，发现电流表示数的变化趋势不同。经检查，发现其中一个实验组使用的滑动变阻器发生断路。

\begin{center}
\begin{tabular}{|c|c|c|c|c|c|}
\hline
\multicolumn{6}{|c|}{表一（第一实验组）} \\ \hline
P 的位置 & a & b & c & d & e \\ \hline
A 的示数（$\UA$） & 0.84 & 0.48 & 0.42 & 0.48 & 0.84 \\ \hline
\end{tabular}

\vspace{0.5em}

\begin{tabular}{|c|c|c|c|c|c|c|}
\hline
\multicolumn{7}{|c|}{表二（第二实验组）} \\ \hline
P 的位置 & a & b & c & d & X & e \\ \hline
A 的示数（$\UA$） & 0.84 & 0.42 & 0.28 & 0.21 & & 0.84 \\ \hline
\end{tabular}
\end{center}

\begin{enumerate}
	\item
	滑动变阻器发生断路的是第\tkanswer[0.8cm]{二}实验组；断路发生在滑动变阻器\tkanswer[1.2cm]{$d\sim e$}段。
	\item
	表二中，对应滑片 P 在 X（d、e 之间的某一点）处的电流表示数的可能值为\xzanswer{D}

	（A）$0.16\UA$\qquad（B）$0.26\UA$\qquad（C）$0.36\UA$\qquad（D）$0.46\UA$
\end{enumerate}

\onepicture[w=4.6cm, align=right]{\includesvg{figs/17}}

%% answer:
\jdanswer{
\begin{enumerate}
	\item
	二；$d\sim e$。
	\item
	D。
\end{enumerate}
}
%% memo:
\memoanswer{
（1）由表一、二中的数据变化情况，可知表一符合并联电路的电阻、电流变化规律，而表二中数据反映电路的电阻在点 d 之前总是增大，故可判断断路发生在第二实验组，且在 $d\sim e$ 之间的某点。

（2）断路在 $d\sim e$ 之间的某点，则滑片 P 可能接触在断点上方，也可能接触在断点下方，如图所示。

\onepicture[w=4.6cm]{figs/17_an.jpg}

设电源电动势为 $E$，电流表 A、电阻 $R$、电源内阻 $r$ 之和为 $R_{\text{外}}$，滑动变阻器的电阻为 $R$，由第一组实验知，当 P 的位置在 a 时，滑动变阻器接入电路的电阻为 0，由闭合电路的欧姆定律得 $\frac{E}{R_{\text{外}}}=0.84\UA$；当 P 的位置在 c 时，滑动变阻器接入电路的总电阻为 $\frac{R}{4}$，由闭合电路欧姆定律得 $\frac{E}{R_{\text{外}}+\frac{R}{4}}=0.42\UA$，联立解得 $R_{\text{外}}=\frac{R}{4}$，即 $R=4R_{\text{外}}$。

当 P 点接触在断点上方时，若 P 点接在 e 时，滑动变阻器接入电路的电阻为 0，此时电流表的示数为 $0.84\UA$；若 P 点接在 d 时，假设端点也恰好在 d 点，此时滑动变阻器接入电路的电阻为 $\frac{R}{4}$，此时电流表的示数为 $0.42\UA$，故电流表的示数介于 $0.42\UA\sim0.84\UA$ 之间，故 D 正确。

当 P 点接触在断点下方时，此时滑动变阻器接入电路的电阻在 $\frac{3R}{4}\sim R$ 之间，当接入电路的电阻为 $\frac{3R}{4}$ 时，电流表的示数为 $I_{\max}=\frac{E}{R_{\text{外}}+\frac{3R}{4}}=\frac{4E}{4R_{\text{外}}+3R}=\frac{E}{4R_{\text{外}}}=0.21\UA$；当接入电路的电阻为 $R$ 时，电流表的示数为 $I_{\min}=\frac{E}{R_{\text{外}}+R}=\frac{E}{5R_{\text{外}}}=0.168\UA$，故电流表示数在 $0.168\UA\sim0.21\UA$ 间，故 ABC 错误。
}

\item
%% number: 18
%% paperName: 2005年上海高考第 18 题 7 分
%% typeId: 4
%% type: 实验题
%% chapter: 运动和力
%% point: 变加速直线运动
%% method: 图象法
%% score: 7
%% degree: 450
%% duplicateId: 0
%% body:
科学探究活动通常包括以下环节：提出问题，作出假设，制定计划，搜集证据，评估交流等。一组同学研究“运动物体所受空气阻力与运动速度关系”的探究过程如下：

A．有同学认为：运动物体所受空气阻力可能与其运动速度有关。

B．他们计划利用一些“小纸杯”作为研究对象，用超声测距仪等仪器测量“小纸杯”在空中直线下落时的下落距离、速度随时间变化的规律，以验证假设。

C．在相同的实验条件下，同学们首先测量了单只“小纸杯”在空中下落过程中不同时刻的下落距离，将数据填入下表中，图\ref{2005上海18a}是对应的位移—时间图线。然后将不同数量的“小纸杯”叠放在一起从空中下落，分别测出它们的速度—时间图线，如图\ref{2005上海18b}中图线 1、2、3、4、5 所示。

\begin{center}
\begin{tabular}{|c|c|}
\hline
时间（$\Us$） & 下落距离（$\Um$） \\ \hline
0.0 & 0.000 \\ \hline
0.4 & 0.036 \\ \hline
0.8 & 0.469 \\ \hline
1.2 & 0.957 \\ \hline
1.6 & 1.447 \\ \hline
2.0 & X \\ \hline
\end{tabular}
\end{center}

D．同学们对实验数据进行分析、归纳后，证实了他们的假设。

回答下列提问：

\begin{enumerate}
	\item
	与上述过程中 A、C 步骤相应的科学探究环节分别是\tkanswer[1.8cm]{作出假设}、\tkanswer[1.8cm]{搜集证据}。
	\item
	图\ref{2005上海18a}中的 AB 段反映了运动物体在做\tkanswer[1.6cm]{匀速}运动，表中 X 处的值为\tkanswer[1.6cm]{1.937}。
	\item
	图\ref{2005上海18b}中各条图线具有共同特点，“小纸杯”在下落的开始阶段做\tkanswer[2.4cm]{加速度逐渐减小的加速}运动，最后“小纸杯”做\tkanswer[1.2cm]{匀速}运动。
	\item
	比较图\ref{2005上海18b}中的图线 1 和 5，指出在 $1.0\Us\sim1.5\Us$ 时间段内，速度随时间变化关系的差异：\tkanswer[5.0cm]{图线 1 反映速度不随时间变化，图线 5 反映速度随时间继续增大}。
\end{enumerate}

\twopicture[wA=6.2cm, wB=6.2cm, labelA=2005上海18a, labelB=2005上海18b]{figs/18a.png}{figs/18b.png}

%% answer:
\jdanswer{
\begin{enumerate}
	\item
	作出假设、搜集证据。
	\item
	匀速运动，1.937。
	\item
	加速度逐渐减小的加速运动，匀速运动。
	\item
	图线 1 反映速度不随时间变化，图线 5 反映速度随时间继续增大（或图线 1 反映纸杯做匀速运动，图线 5 反映纸杯依然在做加速度减小的加速运动）。
\end{enumerate}
}
%% memo:
\memoanswer{
（1）作出假设；搜集证据。

（2）匀速；$1.937$。

（3）加速度逐渐减小的加速；匀速。

（4）图线 1 反映速度不随时间变化，图线 5 反映速度随时间继续增大（或图线 1 反映纸杯做匀速运动，图线 5 反映纸杯依然在做加速度减小的加速运动）。
}

\item
%% number: 19
%% paperName: 2005年上海高考第 19 题 10 分
%% typeId: 6
%% type: 计算题
%% chapter: 运动和力
%% point: 牛顿定律的应用
%% method: 无
%% score: 10
%% degree: 750
%% duplicateId: 0
%% body:
\textbf{A．}某滑板爱好者在离地 $h=1.8\Um$ 高的平台上滑行，水平离开 A 点后落在水平地面的 B 点，其水平位移 $s_{1}=3\Um$，着地时由于存在能量损失，着地后速度变为 $v=4\Ums$，并以此为初速沿水平地面滑行 $s_{2}=8\Um$ 后停止。已知人与滑板的总质量 $m=60\Ukg$。求：

\begin{enumerate}
	\item
	人与滑板在水平地面滑行时受到的平均阻力大小；
	\item
	人与滑板离开平台时的水平初速度。（空气阻力忽略不计，$g=10\Umsq$）
\end{enumerate}

\onepicture[w=6.6cm, align=right]{\includesvg{figs/19a}}

\textbf{B．}如图所示，某人乘雪橇从雪坡经 A 点滑至 B 点，接着沿水平路面滑至 C 点停止。人与雪橇的总质量为 $70\Ukg$。表中记录了沿坡滑下过程中的有关数据，请根据图表中的数据解决下列问题：

\begin{center}
\begin{tabular}{|c|c|c|c|}
\hline
位置 & A & B & C \\ \hline
速度（$\Ums$） & 2.0 & 12.0 & 0 \\ \hline
时刻（$\Us$） & 0 & 4 & 10 \\ \hline
\end{tabular}
\end{center}

\begin{enumerate}
	\item
	人与雪橇从 A 到 B 的过程中，损失的机械能为多少？
	\item
	设人与雪橇在 BC 段所受阻力恒定，求阻力大小。（$g=10\Umsq$）
\end{enumerate}

\onepicture[w=6.6cm, align=right]{\includesvg{figs/19b}}

%% answer:
\jdanswer{
\textbf{A．}
\begin{enumerate}
	\item
	$f=60\UN$。
	\item
	$v_{0}=5\Ums$。
\end{enumerate}

\textbf{B．}
\begin{enumerate}
	\item
	$\Delta E=9100\UJ$。
	\item
	$f=140\UN$。
\end{enumerate}
}
%% memo:
\memoanswer{
\textbf{A．}（1）设滑板在水平地面滑行时受到的平均阻力为 $f$，根据动能定理有

$-fs_{2}=0-\frac{1}{2}mv^{2}$ ①

由①式解得 $f=\frac{mv^{2}}{2s_{2}}=\frac{60\times4^{2}}{2\times8}\UN=60\UN$ ②

（2）人和滑板水平离开 A 点后一起在空中做平抛运动，设初速为 $v_{0}$，飞行时间为 $t$，根据平抛运动规律有

$h=\frac{1}{2}gt^{2}$ ③

$s_{1}=v_{0}t$ ④

由③、④两式解得 $v_{0}=\frac{s_{1}}{\sqrt{\frac{2h}{g}}}=\frac{3}{\sqrt{\frac{2\times1.8}{10}}}\Ums=5\Ums$ ⑤

\textbf{B．}（1）从 A 到 B 的过程中，人与雪橇损失的机械能为

$\Delta E=mgh+\frac{1}{2}mv_{\mathrm{A}}^{2}-\frac{1}{2}mv_{\mathrm{B}}^{2}$ ①

解得 $\Delta E=(70\times10\times20+0.5\times70\times2.0^{2}-0.5\times70\times12.0^{2})\UJ=9100\UJ$ ②

（2）人与雪橇在 BC 段做减速运动的加速度

$a=\frac{v_{\mathrm{C}}-v_{\mathrm{B}}}{t}=\frac{0-12}{10-4}\Umsq=-2\Umsq$ ③

根据牛顿第二定律

$f=ma=70\times(-2)\UN=-140\UN$ ④
}

\item
%% number: 20
%% paperName: 2005年上海高考第 20 题 10 分
%% typeId: 5
%% type: 辨析题
%% chapter: 电场
%% point: 带电粒子的直线运动
%% method: 无
%% score: 10
%% degree: 650
%% duplicateId: 0
%% body:
如图所示，带正电小球质量为 $m=1\times10^{-2}\Ukg$，带电量为 $q=1\times10^{-6}\UC$，置于光滑绝缘水平面上的 A 点。当空间存在着斜向上的匀强电场时，该小球从静止开始始终沿水平面做匀加速直线运动，当运动到 B 点时，测得其速度 $v_{\mathrm{B}}=1.5\Ums$，此时小球的位移为 $s=0.15\Um$。求此匀强电场场强 $E$ 的取值范围。（$g=10\Umsq$）

某同学求解如下：设电场方向与水平面之间夹角为 $\theta$，由动能定理 $qEs\cos\theta=\frac{1}{2}mv_{\mathrm{B}}^{2}-0$ 得 $E=\frac{mv_{\mathrm{B}}^{2}}{2qs\cos\theta}=\frac{75000}{\cos\theta}$。由题意可知 $\theta>0$，所以当 $E>7.5\times10^{4}\UVm$ 时小球将始终沿水平面做匀加速直线运动。经检查，计算无误。该同学所得结论是否有不完善之处？若有请予以补充。

\onepicture[w=5.6cm, align=right]{\includesvg{figs/20}}

%% answer:
\jdanswer{该同学所得结论有不完善之处。$7.5\times10^{4}\UVm<E\le1.25\times10^{5}\UVm$。}
%% memo:
\memoanswer{
该同学所得结论有不完善之处。

为使小球始终沿水平面运动，电场力在竖直方向的分力必须小于等于重力，即 $qE\sin\theta\le mg$ ①，解得 $E\le\frac{mg}{q\sin\theta}$。

水平方向有 $qEs\cos\theta=\frac{1}{2}mv_{\mathrm{B}}^{2}-0$，解得 $qE\cos\theta=\frac{mv_{\mathrm{B}}^{2}}{2s}$ ②。

由 $\frac{\text{①}}{\text{②}}$ 知 $\tan\theta\le\frac{mg}{\frac{mv_{\mathrm{B}}^{2}}{2s}}=\frac{2gs}{v_{\mathrm{B}}^{2}}=\frac{4}{3}$ ③。

由③知 $\sin\theta$ 的最大值取 $\frac{4}{5}$，代入数据得 $E\le\frac{mg}{q\sin\theta}=1.25\times10^{5}\UVm$。又 $\theta>0$，由②得 $E>7.5\times10^{4}\UVm$。综上可知，匀强电场场强 $E$ 的取值范围为 $7.5\times10^{4}\UVm<E\le1.25\times10^{5}\UVm$。
}

\item
%% number: 21
%% paperName: 2005年上海高考第 21 题 10 分
%% typeId: 6
%% type: 计算题
%% chapter: 气体性质
%% point: 理想气体状态方程
%% method: 无
%% score: 10
%% degree: 650
%% duplicateId: 0
%% body:
内壁光滑的导热气缸竖直浸放在盛有冰水混合物的水槽中，用不计质量的活塞封闭压强为 $1.0\times10^{5}\UPa$、体积为 $2.0\times10^{-3}\Umc$ 的理想气体。现在活塞上方缓缓倒上沙子，使封闭气体的体积变为原来的一半，然后将气缸移出水槽，缓慢加热，使气体温度变为 $127\UduC$。

\begin{enumerate}
	\item
	求气缸内气体的最终体积；
	\item
	在 $p-V$ 图上画出整个过程中气缸内气体的状态变化。（大气压强为 $1.0\times10^{5}\UPa$）
\end{enumerate}

\onepicture[w=5.6cm, align=right]{\includesvg{figs/21}}

%% answer:
\jdanswer{
\begin{enumerate}
	\item
	$V_{2}\approx1.47\times10^{-3}\Umc$。
	\item
	如图所示。
\end{enumerate}
}
%% memo:
\memoanswer{
（1）在活塞上方倒沙的过程中温度保持不变，由玻意耳定律得

$p_{0}V_{0}=p_{1}V_{1}$ ①

由①式解得 $p_{1}=\frac{V_{0}}{V_{1}}p_{0}=2\times10^{5}\UPa$ ②

在缓慢加热到 $127\UduC$ 的过程中压强保持不变，由盖-吕萨克定律得

$\frac{V_{1}}{T_{0}}=\frac{V_{2}}{T_{2}}$ ③

由③式解得 $V_{2}=\frac{T_{2}}{T_{0}}V_{1}=\frac{273+127}{273}\times1.0\times10^{-3}\Umc\approx1.47\times10^{-3}\Umc$ ④

（2）在 $p-V$ 图上画出整个过程中气缸内气体的状态变化，如图所示。

\onepicture[w=5.6cm]{\includesvg{figs/21_an}}
}

\item
%% number: 22
%% paperName: 2005年上海高考第 22 题 14 分
%% typeId: 6
%% type: 计算题
%% chapter: 电磁感应
%% point: 电磁感应中的变加速
%% method: 无
%% score: 14
%% degree: 650
%% duplicateId: 0
%% body:
如图所示，处于匀强磁场中的两根足够长、电阻不计的平行金属导轨相距 $1\Um$，导轨平面与水平面成 $\theta=37^{\circ}$ 角，下端连接阻值为 $R$ 的电阻。匀强磁场方向与导轨平面垂直。质量为 $0.2\Ukg$、电阻不计的金属棒放在两导轨上，棒与导轨垂直并保持良好接触，它们之间的动摩擦因数为 0.25。

\begin{enumerate}
	\item
	求金属棒沿导轨由静止开始下滑时的加速度大小；
	\item
	当金属棒下滑速度达到稳定时，电阻 $R$ 消耗的功率为 $8\UW$，求该速度的大小；
	\item
	在上问中，若 $R=2\UO$，金属棒中的电流方向由 a 到 b，求磁感应强度的大小与方向。（$g=10\Umsq$，$\sin37^{\circ}=0.6$，$\cos37^{\circ}=0.8$）
\end{enumerate}

\onepicture[w=5.0cm, align=right]{\includesvg{figs/22}}

%% answer:
\jdanswer{
\begin{enumerate}
	\item
	$a=4\Umsq$。
	\item
	$v=10\Ums$。
	\item
	$B=0.4\UT$，磁场方向垂直导轨平面向上。
\end{enumerate}
}
%% memo:
\memoanswer{
（1）金属棒开始下滑的初速为零，根据牛顿第二定律

$mg\sin\theta-\mu mg\cos\theta=ma$ ①

由①式解得 $a=10\times(0.6-0.25\times0.8)\Umsq=4\Umsq$ ②

（2）设金属棒运动达到稳定时，速度为 $v$，所受安培力为 $F$，棒在沿导轨方向受力平衡

$mg\sin\theta-\mu mg\cos\theta-F=0$ ③

此时金属棒克服安培力做功的功率等于电路中电阻 $R$ 消耗的电功率

$Fv=P$ ④

由③、④两式解得 $v=\frac{P}{F}=\frac{8}{0.2\times10\times(0.6-0.25\times0.8)}\Ums=10\Ums$ ⑤

（3）设电路中电流为 $I$，两导轨间金属棒的长为 $l$，磁场的磁感应强度为 $B$

$I=\frac{Blv}{R}$ ⑥

$P=I^{2}R$ ⑦

由⑥、⑦两式解得 $B=\frac{\sqrt{PR}}{vl}=\frac{\sqrt{8\times2}}{10\times1}\UT=0.4\UT$ ⑧

磁场方向垂直导轨平面向上。
}

\item
%% number: 23
%% paperName: 2005年上海高考第 23 题 14 分
%% typeId: 6
%% type: 计算题
%% chapter: 电路
%% point: 传感器
%% method: 图象法
%% score: 14
%% degree: 250
%% duplicateId: 0
%% body:
一水平放置的圆盘绕竖直固定轴转动，在圆盘上沿半径开有一条宽度为 $2\Umm$ 的均匀狭缝。将激光器与传感器上下对准，使二者间连线与转轴平行，分别置于圆盘的上下两侧，且可以同步地沿圆盘半径方向匀速移动，激光器连续向下发射激光束。在圆盘转动过程中，当狭缝经过激光器与传感器之间时，传感器接收到一个激光信号，并将其输入计算机，经处理后画出相应图线。图\ref{2005上海23a}为该装置示意图，图\ref{2005上海23b}为所接收的光信号随时间变化的图线，横坐标表示时间，纵坐标表示接收到的激光信号强度，图中 $\Delta t_{1}=1.0\times10^{-3}\Us$，$\Delta t_{2}=0.8\times10^{-3}\Us$。

\twopicture[wA=6.4cm, wB=7.2cm, labelA=2005上海23a, labelB=2005上海23b]{figs/23a.png}{figs/23b.png}

\begin{enumerate}
	\item
	利用图\ref{2005上海23b}中的数据求 $1\Us$ 时圆盘转动的角速度；
	\item
	说明激光器和传感器沿半径移动的方向；
	\item
	求图\ref{2005上海23b}中第三个激光信号的宽度 $\Delta t_{3}$。
\end{enumerate}

%% answer:
\jdanswer{
\begin{enumerate}
	\item
	$\omega=7.85\Urads$。
	\item
	激光器和探测器沿半径由中心向边缘移动。
	\item
	$0.67\times10^{-3}\Us$。
\end{enumerate}
}
%% memo:
\memoanswer{
（1）由图线读得，转盘的转动周期 $T=0.8\Us$ ①

角速度 $\omega=\frac{2\pi}{T}=\frac{6.28}{0.8}\Urads=7.85\Urads$ ②

（2）激光器和探测器沿半径由中心向边缘移动（理由为：由于脉冲宽度在逐渐变窄，表明光信号能通过狭缝的时间逐渐减少，即圆盘上对应探测器所在位置的线速度逐渐增加，因此激光器和探测器沿半径由中心向边缘移动）。

（3）设狭缝宽度为 $d$，探测器接收到第 $i$ 个脉冲时距转轴的距离为 $r_{i}$，第 $i$ 个脉冲的宽度为 $\Delta t_{i}$，激光器和探测器沿半径的运动速度为 $v$。

$\Delta t_{i}=\frac{d}{2\pi r_{i}}T$ ③

$r_{3}-r_{2}=r_{2}-r_{1}=vT$ ④

$r_{2}-r_{1}=\frac{dT}{2\pi}\left(\frac{1}{\Delta t_{2}}-\frac{1}{\Delta t_{1}}\right)$ ⑤

$r_{3}-r_{2}=\frac{dT}{2\pi}\left(\frac{1}{\Delta t_{3}}-\frac{1}{\Delta t_{2}}\right)$ ⑥

由④、⑤、⑥式解得 $\Delta t_{3}=\frac{\Delta t_{1}\Delta t_{2}}{2\Delta t_{1}-\Delta t_{2}}=\frac{1.0\times10^{-3}\times0.8\times10^{-3}}{2\times1.0\times10^{-3}-0.8\times10^{-3}}\Us\approx0.67\times10^{-3}\Us$。
}

\end{enumerate}

\end{document}
